% Page 075

\seite{75}

\margmark{15. Februar.}%
    Heute fand die Neuwahl des Gemeinderates statt.\ob
Bei diesem Wahlvorgange des Ortes ist es gelun\ohb
gen, daß in dieser Wahl ein einheitlicher Wahlvorschlag\ob
aufgestellt wurde. Der neue Gemeinderat setzt sich wie\ob
folgt zusammen:

\pend
\abschnitt{1. Schmitz Gaspar}
\pstart


\pend
\begin{verse}
2.\ Sasbach Heinrich\\
3.\ Krämer Nikolaus\\
4.\ Wetraudy-Schirz Matth.\\
5.\ Klassen Wilhelm\\
6.\ Hoffmann Jakob.
\end{verse}
\pstart

\margmark{30. März}%
Am Freitag den 26.\ d. Mts. fand eine Revision der\ob
Schule durch Herrn Kreisschulrat Marbach statt. (2--6 Uhr)\ob
Daran schloß sich um 4--6 Uhr die Abschlußprüfung\ob
der Fortbildungsschule, zu welcher außer dem Herrn\ob
Kreisschulrat auch Herr Bürgermeister Wirtz erschienen\ob
war. Die Bitburger Zeitung vom 28.\ 3. 25 äußerte sich\ob
darüber wie folgt:


\pend
\begin{verse}
Am letzten Freitag nachmittag hielt
\end{verse}
\pstart
Herr Kreisschulrat Marbach die 1.\ Abschlußprüfung\ob
der ländlichen Fortbildungsschule in Sefferweich ab,\ob
zu der auch der Herr Bürgermeister Wirtz erschie\ohb
nen war. Aus dem vorliegenden Gesamtbilde der bisherigen\ob
Unterrichtsarbeit wurde ersichtlich, daß die Grundsätze\ob
beim Unterricht von praktischen Gesichtspunkten aus\ohb
gehen. Die Schüler wurden in Sprachfertig\ohb
keit, Lesen u.\ Landwirtschaftswissen, sodann\ob
über Rechnen und Raumlehre und Zeichenunterricht geprüft. Die\ob
erzielten Ergebnisse erwiesen sich dem Herrn Kreisschulrat\ob
als befriedigend und es\unsicher{} Lehren und Wollen\unsicher\ob
und Sorge geleistet werden. Da die schriftlichen\ob
Ausarbeitungen und damit\unsicher{} die\unsicher{} worden in\ob
Kraft. Der Herr Bürgermeister dankte dem Herrn Kreisschul\ohb
rat, den Schülern für die aufgewandte und gewissen\ohb
hafte Arbeit, die sie nicht nur in ihrem eigenen\ob
Interesse, sondern auch in dem der Gemeinde\ob
und des Vaterlandes geleistet hätten und sprach
\pend
